% Nouveau EGA V (cote 25): archive cote 25, site page 141 (PDF page 142).
% typescript. Draft compared: none.
% AI-transcribed by Claude Opus 5.5 (claude-opus-5-5, adaptive thinking, effort high), Message Batches API, batch ega25b (msgbatch_01A7n5YuXubUPNK1zsepi8RG), 2026-11-13; not yet checked by eye.
\scanpage{25/141}
\mspage{\hand{2}}

\uline{Remarques }4. \var{\inserted[typed]{a)}}{a)}{insertion}{ega25_25_141-v0001} Du \var{\xout{xx}}{}{typewriter strike-out}{ega25_25_141-v0002} critère qui sera développé tout de suite, il résul-%
te que dans la définition 3, il suffit de tester avec un seul point
à valeurs dans un k' - on peut\var{\xout{s}}{}{typewriter strike-out}{ega25_25_141-v0003} par exemple \unsure{prendre|pirendre} $k'=k(x)$ ou $\hand{\overline{\textcolor{black}{k(x)}}}$, et
le point canonique à valeurs dans ce k'. b) Il résulte de la remar
que 2 que x est géométriquement non singulier pour $\hand{\phi}$ sss $\hand{\phi}_x \hand{\neq} 0$ et
$V(\hand{\phi})$ est \uline{lisse} sur k en x. c) Supposons que l'on ait \var{\inserted[typed]{un}}{un}{insertion}{ega25_25_141-v0004} préschéma X lisse
sur un autre \unsure{S|Y}, une section $\hand{\phi}$ de $\underline{O}_X$ et un\var{\xout{e}}{}{typewriter strike-out}{ega25_25_141-v0005} $x \hand{\in} V(\hand{\phi})$, alors on dit que x
est un zéro singulier (resp. supersingulier) \uline{rel à \unsure{S|Y}}, s'il l'est
rel à k(s) sur la fibre $X_s$ (s étant l'image de x dans \unsure{S|Y}). d) Sous
les conditions de déf 1, on voit tout de suite que la singularité resp.
\var{\margin{$\|$}}{\placedmargin{$\|$}}{marginal note placed}{ega25_25_141-v0006} supersingularité d'un $x \hand{\in} V(\hand{\phi})$ pour $\hand{\phi}$ n'est pas modifiée, quand on
remplace $\hand{\phi}$ par $\hand{\phi}'=u\hand{\phi}$ , où u est une unité en x. Il s'ensuit aussitôt que
la \var{\xout{notion}}{}{typewriter strike-out}{ega25_25_141-v0007} définition 1 et par suite aussi la déf. 3 s'étendent de façon
évidente au cas où $\hand{\phi}$ est une section d'un Module inversible $\underline{L}$ (de façon
à redonner les définitions initiales pour $L = \underline{O}_X$) .

\var{\xout{Considérons la section xxxxx d$^2$}}{}{typewriter strike-out}{ega25_25_141-v0008} Soit X un préschéma lisse sur un
autre Y, et soit $\hand{\phi}$ une section de $\underline{O}_X$ , d'où une section $d^2_{X/Y}\hand{\phi}$ de
$P^2_{X/Y}$ , se réduisant en la section $d^1_{X/Y}\hand{\phi}$ de $P^1_{X/Y}$ , laquelle elle-même
% NOTE: from "Ceci posé" to the end of the page the text is enclosed by hand in a frame, with small tick marks along its lines; the marginal note on the left refers to it
se réduit en la section $d^0\hand{\phi}$ \var{\inserted[typed]{$=\hand{\phi}$}}{$=\hand{\phi}$}{insertion}{ega25_25_141-v0009} de $P^0_{X/Y} = \underline{O}_X$ . Ceci posé, on a la
\uline{Proposition} 5 L'ensemble des zéros de $d^{\hand{0}}\hand{\phi}$ ( resp $d^1\hand{\phi}$ \hand{)} \var{\deleted{resp. $d^2\hand{\phi}$)}}{}{deletion}{ega25_25_141-v0010}
\var{\margin{garder le texte ci-contre}}{\placedmargin{garder le texte ci-contre}}{marginal note placed}{ega25_25_141-v0011} est respectivement égal à l'ensemble $V(\hand{\phi})$ des zéros de $\hand{\phi}$ (resp. à
l'ensemble $V(\hand{\phi})^{\mathrm{sing}}$ des zéros de $\hand{\phi}$ singuliers rel à S\hand{)} \var{\deleted{resp. à l'ens.}}{}{deletion}{ega25_25_141-v0012}
\var{\deleted{$V(\hand{\phi})^{\mathrm{sup\,sing}}$ des zéros de $\hand{\phi}$ supersinguliers rel à S)}}{}{deletion}{ega25_25_141-v0013}.

La première assertion est triviale, la deuxième n'est autre
que le critère jacobien, ou si on préfère résulte de l'isomorphisme
canonique $\underline{m}_x/\underline{m}_x^2 \hand{\simeq} \underline{\hand{\Omega}}^1_{X/k}(x)$ lorsque x est un point rat/k d'un

\mspage{\hand{141}}
